\section{Các công thức lượng giác}

\subsection{Đề 1}
\Opensolutionfile{ans}[ans/ans-1-C3B3-DE1-LC]
\caulc
%%==========Câu 1
\begin{ex}%[1D1N3-2]
Tính $\sin{105^\circ}$ ta được
\choice
{$\dfrac{\sqrt{6}-\sqrt{2}}{4}$}
{$-\dfrac{\sqrt{6}-\sqrt{2}}{4}$}
{\True $\dfrac{\sqrt{6}+\sqrt{2}}{4}$}
{$-\dfrac{\sqrt{6}+\sqrt{2}}{4}$}
\loigiai{
	Có $\sin{105^\circ}=\sin (60^\circ+45^\circ)$ $=$ $\sin{60^\circ}\cdot \cos{45^\circ}+\cos{60^\circ}\cdot \sin{45^\circ}$.\\
	$\Rightarrow$ $\sin{105^\circ}=\dfrac{\sqrt{3}}{2}\cdot \dfrac{\sqrt{2}}{2}+\dfrac{1}{2}\cdot \dfrac{\sqrt{2}}{2}$ $=$ $\dfrac{\sqrt{6}+\sqrt{2}}{4}$.}
\end{ex}

%%==========Câu 2
\begin{ex}%[1D1H3-2]
Rút gọn biểu thức $\sin (a-17^\circ)\cdot \cos (a+13^\circ)-\sin (a+13^\circ)\cdot \cos (a-17^\circ)$,
ta được
\choice
{$\sin 2a$}
{$\cos 2a$}
{\True $-\dfrac{1}{2}$}
{$\dfrac{1}{2}$}
\loigiai{
	Ta có $\sin (a-17^\circ)\cdot \cos (a+13^\circ)-\sin (a+13^\circ)\cdot \cos (a-17^\circ)=\sin \left[(a-17^\circ)-(a+13^\circ)\right]$
	$=\sin (-30^\circ)=-\dfrac{1}{2}$.}
\end{ex}

%%==========Câu 3
\begin{ex}%[1D1H3-2]
Nếu biết $\sin a=\dfrac{8}{17}$, $\tan b=\dfrac{5}{12}$ và $a$, $b$ đều là các góc nhọn và dương thì $\sin (a-b)$ là
\choice
{$\dfrac{20}{220}$}
{$-\dfrac{20}{220}$}
{\True $\dfrac{21}{221}$}
{$\dfrac{22}{221}$}
\loigiai{
	Ta có $a$, $b$ đều là các góc nhọn và dương.\\
	$\sin a=\dfrac{8}{17}\Rightarrow \cos a=\sqrt{1-\dfrac{64}{289}}=\dfrac{15}{17}$.\\
	$\tan b=\dfrac{5}{12}\Rightarrow \cos b=\dfrac{1}{\sqrt{1+\dfrac{25}{144}}}=\dfrac{12}{13}\Rightarrow \sin b=\tan b\cdot \cos b=\dfrac{5}{13}$.\\
	$\Rightarrow \sin (a-b)=\dfrac{8}{17}\cdot \dfrac{12}{13}-\dfrac{15}{17}\cdot \dfrac{5}{13}=\dfrac{21}{221}$.}
\end{ex}

%%==========Câu 4
\begin{ex}%[1D1H3-2]
Gọi $M=\cos (a+b)\cdot \cos (a-b)-\sin (a+b)\cdot \sin (a-b)$ thì
\choice
{$M=1-2\cos^2a$}
{\True $M=1-2\sin^2a$}
{$M=\cos 4a$}
{$M=\sin 4a$}
\loigiai{
	Ta có 
	\begin{eqnarray*}
		M&=&\cos (a+b)\cdot \cos (a-b)-\sin (a+b)\cdot \sin (a-b)\\
		&=&	\cos (a+b+a-b)=\cos 2a=1-2\sin^2a.
	\end{eqnarray*}
	}
\end{ex}

%%==========Câu 5
\begin{ex}%[1D1H3-3]
Biểu thức $\dfrac{1+\sin 4\alpha -\cos4\alpha}{1+\sin 4\alpha +\cos4\alpha}$ có kết quả rút gọn bằng
\choice
{$\sin 2\alpha$}
{$\cos 2\alpha$}
{\True $\tan 2\alpha$}
{$\cot 2\alpha$}
\loigiai{
	$\dfrac{1+\sin 4\alpha -\cos4\alpha}{1+\sin 4\alpha +\cos4\alpha}=\dfrac{2\sin^22\alpha +2\sin 2\alpha \cos2\alpha}{2\cos^22\alpha +2\sin 2\alpha \cos2\alpha}=\dfrac{2\sin 2\alpha (\sin 2\alpha +\cos2\alpha)}{2\cos 2\alpha (\sin 2\alpha +\cos2\alpha)}=\tan 2\alpha$.}
\end{ex}

%%==========Câu 6
\begin{ex}%[1D1H3-3]
Nếu $\tan \dfrac{x}{2}=\dfrac{a}{b}$ thì biểu thức $a\sin x+b\cos x$ bằng
\choice
{$a$}
{\True $b$}
{$\dfrac{a+b}{a}$}
{$\dfrac{a+b}{b}$}
\loigiai{
	Đặt $t=\tan \dfrac{x}{2}=\dfrac{a}{b}$ nên 
	\begin{eqnarray*}
	&\sin x=\dfrac{2t}{1+t^2}=\dfrac{2\dfrac{a}{b}}{1+\dfrac{a^2}{b^2}}=\dfrac{2ab}{a^2+b^2}\\
	&\cos x=\dfrac{1-t^2}{1+t^2}=\dfrac{1-\dfrac{a^2}{b^2}}{1+\dfrac{a^2}{b^2}}=\dfrac{b^2-a^2}{a^2+b^2}.
	\end{eqnarray*}
	Vậy $a\sin x+b\cos x=\dfrac{2a^2b}{a^2+b^2}+\dfrac{b^3-a^2b}{a^2+b^2}=b$.}
\end{ex}

%%==========Câu 7
\begin{ex}%[1D1H3-1]
Cho $M=6\cos^2x+5\sin^2x$. Khi đó giá trị lớn nhất của $M$ là
\choice
{$1$}
{$5$}
{\True $6$}
{$11$}
\loigiai{
	$M=6(1-\sin^2x)+5\sin^2x=6-\sin^2x$.\\
	Ta có $0\le\sin^2x\le 1,\forall x\in \mathbb{R}\Leftrightarrow 0\ge -\sin^2x\ge -1,\forall x\in \mathbb{R}\Leftrightarrow 6\ge 6-\sin^2x\ge 5,\forall x\in \mathbb{R}$.\\
	Giá trị lớn nhất là $6$.}
\end{ex}

%%==========Câu 8
\begin{ex}%[1D1V3-6]
Biết $A$, $B$, $C$ là các góc của tam giác $ABC$, khi đó.
\choice
{$\cos \left(\dfrac{A+B}{2}\right)=\cos \dfrac{C}{2}$}
{$\cos \left(\dfrac{A+B}{2}\right)=-\cos \dfrac{C}{2}$}
{\True $\tan \left(\dfrac{A+B}{2}\right)=\cot \dfrac{C}{2}$}
{$\cot \left(\dfrac{A+B}{2}\right)=\cot \dfrac{C}{2}$}
\loigiai{
	Vì $A$, $B$, $C$ là các góc của tam giác $ABC$ nên $A+B+C=180^\circ\Rightarrow C=180^\circ-(A+B)$.\\
	$\Rightarrow \dfrac{C}{2}=90^\circ-\dfrac{A+B}{2}$. Do đó $\dfrac{C}{2}$ và $\dfrac{A+B}{2}$ là $2$ góc phụ nhau.\\
	$\Rightarrow \sin \dfrac{C}{2}=\cos \dfrac{A+B}{2}$; $\cos \dfrac{C}{2}=\sin \dfrac{A+B}{2}$; $\tan \dfrac{C}{2}=\cot \dfrac{A+B}{2}$; $\cot \dfrac{C}{2}=\tan \dfrac{A+B}{2}$.}
\end{ex}

%%==========Câu 9
\begin{ex}%[1D1V3-5]
Nếu biết $\heva{&\tan a+\tan b=2 \\&\tan (a+b)=4}$ thì các giá trị của $\tan a$, $\tan b$ bằng
\choice
{$\dfrac{1}{3}$, $\dfrac{5}{3}$ hoặc ngược lại}
{$\dfrac{1}{2}$, $\dfrac{3}{2}$ hoặc ngược lại}
{$1-\dfrac{\sqrt{3}}{2}$, $1+\dfrac{\sqrt{3}}{2}$ hoặc ngược lại}
{\True $1-\dfrac{\sqrt{2}}{2}$, $1+\dfrac{\sqrt{2}}{2}$ hoặc ngược lại}
\loigiai{
	Ta có $\heva{&\tan a+\tan b=2 \\&\tan (a+b)=4.}$\\
	Từ $\tan (a+b)=4\Leftrightarrow \dfrac{\tan a+\tan b}{1-\tan a\cdot \tan b}=4\Rightarrow 2=4-4\tan a\cdot \tan b\Rightarrow \tan a\cdot \tan b=\dfrac{1}{2}$.\\
	$\Rightarrow \tan a$, $\tan b$ theo thứ tự là nghiệm của phương trình $X^2-2X+\dfrac{1}{2}=0$.\\
	$\Rightarrow \tan a=1-\dfrac{\sqrt{2}}{2}$, $\tan b=1+\dfrac{\sqrt{2}}{2}$ hoặc ngược lại.}
\end{ex}

%%==========Câu 10
\begin{ex}%[1D1H3-5]
Trong các hệ thức sau, hệ thức nào \textbf{sai}?
\choice
{$\dfrac{\cos2x}{1+\sin 2x}=\dfrac{1-\tan x}{1+\tan x}$}
{$4\sin a\cdot \cos a(1-2\sin^2a)=\sin 4a$}
{$\cos 4a = 8\cos{^4}a-8\cos^2a+1$}
{\True $\cos 4a-4\cos 2a+3=8\cos^4a$}
\loigiai{
	$\dfrac{\cos2x}{1+\sin 2x}=\dfrac{\cos^2x-\sin^2x}{(\sin x+\cos x)^2}=\dfrac{(\cos x-\sin x)(\sin x+\cos x)}{(\sin x+\cos x)^2}=\dfrac{\cos x-\sin x}{\sin x+\cos x}=\dfrac{1-\tan x}{1+\tan x}$.\\
	$4\sin a\cdot \cos a(1-2\sin^2a)=2\sin 2a\cdot \cos2a=\sin 4a$.\\
	$\cos4a=2\cos^22a-1= 2(2\cos^2a-1)^2=8\cos^4a-8\cos^2a+1$.\\
	$\cos4a- 4\cos2a+3=2(1-2\sin^2a)^2-1-4(1-2\sin^2a)+3=8\sin^4a$.}
\end{ex}

%%==========Câu 11
\begin{ex}%[1D1H3-3]
Nếu $\sin a-\cos a=\dfrac{1}{5}(135^\circ<a<180^\circ)$ thì giá trị đúng của $\tan 2a$ là
\choice
{$-\dfrac{20}{7}$}
{$\dfrac{20}{7}$}
{\True $\dfrac{24}{7}$}
{$-\dfrac{24}{7}$}
\loigiai{
	$\sin a-\cos a=\dfrac{1}{5}\Rightarrow 1-\sin 2a=\dfrac{1}{25}\Rightarrow \sin 2a=\dfrac{24}{25}\Rightarrow \cos 2a=\sqrt{1-\dfrac{576}{625}}=\dfrac{7}{25}$ $\Rightarrow \tan 2a=\dfrac{24}{7}$.}
\end{ex}

%%==========Câu 12
\begin{ex}%[1D1V3-1]
Giá trị lớn nhất của $M=\sin^4x+\cos^4x$ bằng
\choice
{\True $1$}
{$2$}
{$3$}
{$4$}
\loigiai{
	Ta có $M=1-\dfrac{1}{2}\sin^22x$.\\
	Vì $0\le{\sin^2}x\le 1$
	$\Leftrightarrow -\dfrac{1}{2}\le -\dfrac{1}{2}\sin^22x\le 0$
	$\Leftrightarrow \dfrac{1}{2}\le 1-\dfrac{1}{2}\sin^22x\le 1$.\\
	Nên giá trị lớn nhất là $1$.}
\end{ex}

\Closesolutionfile{ans}
\indapan{6}{ans/ans-1-C3B3-DE1-LC}
\cauds
\Opensolutionfile{ans}[ans/ans-1-C3B3-DE1-DS]

%%==========Câu 1
\begin{ex}%[1D1N3-2]
	Cho biết $\sin \alpha=\dfrac{3}{5}$, $\dfrac{\pi}{2}<\alpha<\pi$. Khi đó
\choiceTF
{\True $\cos \alpha <0$}
{\True $\cos \alpha =-\dfrac{4}{5}$}
{$\tan \alpha =\dfrac{3}{4}$}
{$\tan \left(\alpha +\dfrac{\pi}{3}\right)=\dfrac{48-\sqrt{3}}{11}$}
\loigiai{
	\begin{itemchoice}
		\itemch Đúng. Vì $\dfrac{\pi}{2}<\alpha <\pi$ nên $\cos \alpha <0$.
		\itemch Đúng. Ta có $\cos \alpha =-\sqrt{1-\sin^2\alpha}=-\dfrac{4}{5}$.		
		\itemch Sai. Ta có $\tan \alpha =\dfrac{\sin \alpha}{\cos \alpha}=-\dfrac{3}{4}$.
		\itemch Sai. Ta có $\tan \left(\alpha+\dfrac{\pi}{3}\right)=\dfrac{\tan \alpha+\tan \dfrac{\pi}{3}}{1-\tan \alpha \tan \dfrac{\pi}{3}}=\dfrac{\tan \alpha+\sqrt{3}}{1-\sqrt{3} \tan \alpha}$.\\
		Suy ra $\tan \left(\alpha+\dfrac{\pi}{3}\right)=\dfrac{-\dfrac{3}{4}+\sqrt{3}}{1-\sqrt{3}\left(-\dfrac{3}{4}\right)}=\dfrac{-3+4 \sqrt{3}}{4+3 \sqrt{3}}=\dfrac{48-25 \sqrt{3}}{11}$.
	\end{itemchoice}
}
\end{ex}

%%==========Câu 2
\begin{ex}%[1D1H3-2]
	Cho $\cot x=-\sqrt{3}$, $\dfrac{3 \pi}{2}<x<2 \pi$. Khi đó
	\choiceTF
	{\True $\sin x=-\dfrac{\sqrt{10}}{10}$}
	{$\cos x=\dfrac{\sqrt{3}}{10}$}
	{\True $\sin \left(\dfrac{4\pi}{3}-x\right)=\dfrac{-\sqrt{10}}{5}$}
	{\True $\tan \left(x+\dfrac{\pi}{3}\right)=\dfrac{\sqrt{3}}{3}$}
	\loigiai{
		\begin{itemchoice}
			\itemch Đúng. $\cot x=-\sqrt{3}$, $\dfrac{3\pi}{2}<x<2\pi$. Khi đó $\dfrac{1}{\sin^2x}=1+\cot^2x=1+(-\sqrt{3})^2=10\\\Rightarrow \sin^2x=\dfrac{1}{10}\Rightarrow \sin x=\pm \dfrac{\sqrt{10}}{10}$.\\
			Vì $\dfrac{3\pi}{2}<x<2\pi $ nên $\sin x=-\dfrac{\sqrt{10}}{10}$.
			\itemch Sai. $\cos x=\cot x\cdot \sin x=-\sqrt{3}\cdot \left(-\dfrac{\sqrt{10}}{10}\right)=\dfrac{\sqrt{30}}{10}$.
			\itemch Đúng. 
			\begin{eqnarray*}
				\sin \left(\dfrac{4 \pi}{3}-x\right)&=&\sin \dfrac{4 \pi}{3} \cos x-\cos \dfrac{4 \pi}{3} \sin x\\&=&-\dfrac{\sqrt{3}}{2} \cdot\left(\dfrac{\sqrt{30}}{10}\right)-\dfrac{-1}{2} \cdot \dfrac{-\sqrt{10}}{10}=\dfrac{-\sqrt{10}}{5}.
			\end{eqnarray*}
			\itemch Đúng.
			\begin{eqnarray*}
			\tan x&=&\dfrac{1}{\cot x}=-\dfrac{\sqrt{3}}{3}\tan \left(x+\dfrac{\pi}{3}\right)=\dfrac{\tan x+\tan \dfrac{\pi}{3}}{1-\tan x \cdot \tan \dfrac{\pi}{3}}\\&=&\dfrac{\dfrac{-\sqrt{3}}{3}+\sqrt{3}}{1-\dfrac{-\sqrt{3}}{3} \cdot \sqrt{3}}=\dfrac{\sqrt{3}}{3}.	
			\end{eqnarray*}
		\end{itemchoice}
	}
\end{ex}

%%==========Câu 3
\begin{ex}%[1D1H3-3]
Cho $\cos \alpha=\dfrac{2}{5}$, $2 \pi<\alpha<\dfrac{5 \pi}{2}$, khi đó	
	\choiceTF
	{\True $\sin \alpha=\dfrac{\sqrt{21}}{5}$}
	{\True $\sin 2\alpha =\dfrac{4\sqrt{21}}{25}$}
	{\True $\cos 2\alpha =-\dfrac{17}{25}$}
	{\True $\tan 2\alpha =\dfrac{-4\sqrt{21}}{17}$}
	\loigiai{
		\begin{itemchoice}
			\itemch Đúng. $\cos \alpha=\dfrac{2}{5}$, $2 \pi<\alpha<\dfrac{5 \pi}{2}\Rightarrow\sin ^2 \alpha=1-\cos ^2 \alpha=1-\left(\dfrac{2}{5}\right)^2=\dfrac{21}{25} \\\Rightarrow \sin \alpha= \pm \dfrac{\sqrt{21}}{5}$.\\
	Vì $2 \pi<\alpha<\dfrac{5 \pi}{2}$ nên $\sin \alpha=\dfrac{\sqrt{21}}{5}$.
			\itemch Đúng.  $\sin 2 \alpha=2 \sin \alpha \cos \alpha=2 \cdot \dfrac{\sqrt{21}}{5} \cdot \dfrac{2}{5}=\dfrac{4 \sqrt{21}}{25}$.
			\itemch Đúng. $\cos 2 \alpha=\cos ^2 \alpha-\sin ^2 \alpha=\left(\dfrac{2}{5}\right)^2-\left(\dfrac{\sqrt{21}}{5}\right)^2=-\dfrac{17}{25}$.
			\itemch Đúng. $\tan 2 \alpha=\dfrac{\sin 2 \alpha}{\cos 2 \alpha}=\dfrac{4 \sqrt{21}}{25}\colon \dfrac{-17}{25}=\dfrac{-4 \sqrt{21}}{17}$.
		\end{itemchoice}
	}
\end{ex}

%%==========Câu 4
\begin{ex}%[1D1H3-4]
	Biến đổi được các biểu thức sau về dạng tích số. Khi đó
	\choiceTF
	{$\cos 3x+\cos x=2\cos 2x\cdot \cos 3x$}
	{$\sin 3x+\sin 2x=2\sin 2x\cos \dfrac{x}{2}$}
	{\True $\cos 4 x-\cos x=-2 \sin \dfrac{5 x}{2} \sin \dfrac{3 x}{2}$}
	{\True $\sin 5 x-\sin x=2 \cos 3 x \sin 2 x$}
	\loigiai{
		\begin{itemchoice}
			\itemch Sai. $\cos 3 x+\cos x=2 \cos 2 x \cdot \cos x$.
			\itemch Sai. $\sin 3 x+\sin 2 x=2 \sin \dfrac{5 x}{2} \cos \dfrac{x}{2}$.
			\itemch Đúng. $\cos 4 x-\cos x=-2 \sin \dfrac{5 x}{2} \sin \dfrac{3 x}{2}$.
			\itemch Đúng. $\sin 5 x-\sin x=2 \cos 3 x \sin 2 x$.
		\end{itemchoice}
	}
\end{ex}
\Closesolutionfile{ans}
\indapan{3}{ans/ans-1-C3B3-DE1-DS}

\Opensolutionfile{ans}[ans/ans-1-C3B3-DE1-KQ]
\caukq

%%==========Câu 1
\begin{ex}%[1D1H3-1]
	Cho $\sin \alpha=\dfrac{3}{5}$ và $(90^\circ<\alpha<180^\circ)$. Tính $\cos \alpha$.
	\shortans{$-0{,}8$}
	\loigiai{
	Ta có $\sin ^2 \alpha+\cos ^2 \alpha=1 \Rightarrow \cos ^2 \alpha=1-\sin ^2 \alpha=1-\left(\dfrac{3}{5}\right)^2=\dfrac{16}{25} \Rightarrow \cos \alpha= \pm \dfrac{4}{5}$. $90^\circ<\alpha<180^\circ$ nên $\cos \alpha<0$.\\
	Vậy $\cos \alpha=-\dfrac{4}{5}$.	
	}
\end{ex}

%%==========Câu 2
\begin{ex}%[1D1H3-1]
	Cho $\tan \alpha=2$. Tính $C=\dfrac{\sin \alpha}{\sin ^3 \alpha+2 \cos ^3 \alpha}$.
	\shortans{$1$}
	\loigiai{
	$C=\dfrac{\sin \alpha}{\sin ^3 \alpha+2 \cos ^3 \alpha}=\dfrac{\tan \alpha \cdot \dfrac{1}{\cos ^2 \alpha}}{\tan ^3 \alpha+2}=\dfrac{\tan \alpha(1+\tan ^2 \alpha)}{\tan ^3 \alpha+2}=\dfrac{2(1+2^2)}{2^3+2}=1$.	
	}
\end{ex}

%%==========Câu 3
\begin{ex}%[1D1H3-3]
	Biết $\dfrac{3-4\cos 2a+\cos 4a}{3+4\cos 2a+\cos 4a}=k\tan^4a$, tìm $k$?
	\shortans{$1$}
	\loigiai{
Ta có 
\begin{eqnarray*}
\dfrac{3-4 \cos 2 a+\cos 4 a}{3+4 \cos 2 a+\cos 4 a}&=&\dfrac{2 \cos ^2 2 a-4 \cos 2 a+2}{2 \cos ^2 2 a+4 \cos 2 a+2}=\dfrac{2(\cos ^2 2 a-2 \cos 2 a+1)}{2(\cos ^2 2 a+2 \cos 2 a+1)}\\
&=&\dfrac{(1-\cos 2 a)^2}{(1+\cos 2 a)^2}=\dfrac{\left[1-(1-2 \sin ^2 a)\right]^2}{\left[1+(2 \cos ^2 a-1)\right]^2}=\dfrac{4 \sin ^4 a}{4 \cos ^4 a}=\tan ^4 a.
\end{eqnarray*}
}
\end{ex}

%%==========Câu 4
\begin{ex}%[1D1H3-2]
	Cho $\alpha-\beta=\dfrac{\pi}{3}$. Giá trị của biểu thức $B=(\cos \alpha+\sin \beta)^2+(\cos \beta-\sin \alpha)^2$ bằng $a-\sqrt{b}$. Khi đó $a+b$ bằng
	\shortans{$5$}
	\loigiai{
	Ta có
	\begin{eqnarray*}
	B & =&\cos ^2 \alpha+\sin ^2 \beta+2 \cos \alpha \sin \beta+\cos ^2 \beta+\sin ^2 \alpha-2 \sin \alpha \cos \beta \\
	& =&2+2 \sin \beta \cos \alpha-2 \cos \beta \sin \alpha \\
	& =&2+2 \sin (\beta-\alpha)=2+2 \sin \left(-\dfrac{\pi}{3}\right)=2-\sqrt{3}.
	\end{eqnarray*}
	}
\end{ex}

%%==========Câu 5
\begin{ex}%[1D1H3-7]
	Từ một vị trí $A$, người ta buộc hai sợi cáp $A B$ và $A C$ đến một cái trụ cao $15~m$, được dựng vuông góc với mặt đất, chân trụ ở vị trí $D$. Biết $C D=9~m$ và $A D=12~m$. Tìm góc nhọn $\alpha=\widehat{B A C}$ tạo bởi hai sợi dây cáp đó, đồng thời tính gần đúng $\alpha$ (làm tròn đến hàng phần chục, đơn vị độ).
	\shortans{$14{,}5$}
	\loigiai{
	Ta có
	\begin{eqnarray*}
	\tan \alpha&=&\tan (\widehat{B A D}-\widehat{C A D}) \\
	&=&\dfrac{\tan \widehat{B A D}-\tan \widehat{C A D}}{1+\tan \widehat{B A D} \tan \widehat{C A D}}=\dfrac{\dfrac{15}{12}-\dfrac{9}{12}}{1+\dfrac{15}{12} \cdot \dfrac{9}{12}}=\dfrac{8}{31}.
	\end{eqnarray*}
Vì vậy $\alpha \approx 14{,}5^\circ$.	
	}
\end{ex}

%%==========Câu 6
\begin{ex}%[1D1C3-6]
	Cho tam giác $ABC$ có $AB=c$, $AC=b$, $BC=a$, thỏa mãn $$\dfrac{-1+\cos (\pi+B)}{\sin (\pi-A-C)}=-\dfrac{2 a+c}{\sqrt{4 a^2-c^2}}.$$
	Giá trị của biểu thức $(a-b)^2c^3$ bằng
	\shortans{$0$}
	\loigiai{
	Ta có
	\begin{eqnarray*}
	&& \dfrac{-1+\cos (\pi +B)}{\sin (\pi -A-C)}=-\dfrac{2a+c}{\sqrt{4a^2-c^2}}\Leftrightarrow \dfrac{-1-\cos B}{\sin (\pi -(A+C))}=-\dfrac{2a+c}{\sqrt{4a^2-c^2}} \\
	& \Leftrightarrow& \dfrac{1+\cos B}{\sin B}=\dfrac{2a+c}{\sqrt{4a^2-c^2}}\Leftrightarrow \dfrac{1+\cos B}{\sin B}=\dfrac{2a+c}{\sqrt{4a^2-c^2}}\Leftrightarrow{\left(\dfrac{1+\cos B}{\sin B}\right)^2}=\left(\dfrac{2a+c}{\sqrt{4a^2-c^2}}\right)^2 \\
	& \Leftrightarrow& \dfrac{(1+\cos B)^2}{\sin^2B}=\dfrac{(2a+c)^2}{4a^2-c^2}\Leftrightarrow \dfrac{(1+\cos B)^2}{1-\cos^2B}=\dfrac{(2a+c)^2}{4a^2-c^2}\Leftrightarrow \dfrac{1+\cos B}{1-\cos B}=\dfrac{2a+c}{2a-c} \\
	& \Leftrightarrow& \dfrac{1+\cos B}{1-\cos B}-1=\dfrac{2a+c}{2a-c}-1\Leftrightarrow 2ac\cdot \cos B=c^2\Leftrightarrow a^2+c^2-b^2=c^2\Leftrightarrow a^2=b^2\Leftrightarrow a=b.
	\end{eqnarray*}
Vậy $(a-b)^2c^3=0$.	
	}
\end{ex}
\Closesolutionfile{ans}
\indapan{6}{ans/ans-1-C3B3-DE1-KQ}

\subsection{Đề 2}
\Opensolutionfile{ans}[ans/ans-1-C3B3-DE2-LC]
\caulc
%%==========Câu 1
\begin{ex}%[1D1N3-2]
Tính $\tan 105^\circ$ ta được
\choice
{\True $-(2+\sqrt{3})$}
{$2+\sqrt{3}$}
{$2-\sqrt{3}$}
{$-(2-\sqrt{3})$}
\loigiai{
	$\tan{105^\circ}$ $=$ $\tan (45^\circ+60^\circ)$ $=$ $\dfrac{\tan{45^\circ}+\tan{60^\circ}}{1-\tan{45^\circ}\tan{60^\circ}}$ $=$ $\dfrac{1+\sqrt{3}}{1-\sqrt{3}}$ $=$ $-(2+\sqrt{3})$.}
\end{ex}

%%==========Câu 2
\begin{ex}%[1D1N3-4]
Rút gọn biểu thức $\cos (x+\dfrac{\pi}{4})-\cos (x-\dfrac{\pi}{4})$ ta được
\choice
{$\sqrt{2}\sin x$}
{\True $-\sqrt{2}\sin x$}
{$\sqrt{2}\cos x$}
{$-\sqrt{2}\cos x$}
\loigiai{
	Ta có 
	\begin{eqnarray*}
	\cos (x+\dfrac{\pi}{4})-\cos (x-\dfrac{\pi}{4})&=&-2\sin \left(\dfrac{x+\dfrac{\pi}{4}+x-\dfrac{\pi}{4}}{2}\right)\cdot \sin \left(\dfrac{x+\dfrac{\pi}{4}-x+\dfrac{\pi}{4}}{2}\right)\\&=&-\sqrt{2}\sin x.	
	\end{eqnarray*}
	}
\end{ex}

%%==========Câu 3
\begin{ex}%[1D1H3-2]
Nếu $\tan x=0{,}5$; $\sin y=\dfrac{3}{5}(0<y<90^\circ)$ thì $\tan (x+y)$ bằng
\choice
{\True $2$}
{$3$}
{$4$}
{$5$}
\loigiai{
	$\tan x=0{,}5=\dfrac{1}{2}$, $\sin y=\dfrac{3}{5}(0<y<90^\circ)\Rightarrow \cos y=\dfrac{4}{5}\Rightarrow \tan y=\dfrac{3}{4}$.\\
	$\tan (x+y)=\dfrac{\tan x+\tan y}{1-\tan x\cdot \tan y}=\dfrac{\dfrac{1}{2}+\dfrac{3}{4}}{1-\dfrac{1}{2}\cdot \dfrac{3}{4}}=2$.}
\end{ex}

%%==========Câu 4
\begin{ex}%[1D1N3-1]
Gọi $M=\tan x-\tan y$ thì
\choice
{$M=\tan (x-y)$}
{$M=\dfrac{\sin (x+y)}{\cos x\cdot \cos y}$}
{\True $M=\dfrac{\sin (x-y)}{\cos x\cdot \cos y}$}
{$M=\dfrac{\tan x-\tan y}{1+\tan x\cdot \tan y}$}
\loigiai{
	Ta có $M=\tan x-\tan y=\dfrac{\sin x}{\cos x}-\dfrac{\sin y}{\cos y}=\dfrac{\sin x\cos y-\cos x\sin y}{\cos x\cos y}$ $=\dfrac{\sin (x-y)}{\cos x\cos y}$.}
\end{ex}

%%==========Câu 5
\begin{ex}%[1D1H3-1]
Biểu thức $\dfrac{3-4\cos 2\alpha +\cos4\alpha}{3+4\cos 2\alpha +\cos4\alpha}$ có kết quả rút gọn bằng
\choice
{$-\tan^4\alpha$}
{\True $\tan^4\alpha$}
{$-\cot^4\alpha$}
{$\cot^4\alpha$}
\loigiai{
	Ta có 
	\begin{eqnarray*}
	&& \dfrac{3-4\cos 2\alpha +\cos4\alpha}{3+4\cos 2\alpha +\cos4\alpha}=\dfrac{3-4(1-2\sin^2\alpha)+2(1-2\sin^2\alpha)^2-1}{3+4(2\cos^2\alpha -1)+2(2\cos^2\alpha -1)^2-1} \\
	& =&\dfrac{8\sin^2a-8\sin^2\alpha +8\sin^4\alpha}{8\cos^2a-8\cos^2\alpha +8\cos^4\alpha}=\tan^4\alpha.	
	\end{eqnarray*}
	}
\end{ex}

%%==========Câu 6
\begin{ex}%[1D1H3-3]
Cho biểu thức $A=\sin^2(a+b)-\sin^2a-\sin^2b$. Hãy chọn kết quả đúng
\choice
{$A=2\cos a\cdot \sin b\cdot \sin (a+b)$}
{$A=2\sin a\cdot \cos b\cdot \cos (a+b)$}
{$A=2\cos a\cdot \cos b\cdot \cos (a+b)$}
{\True $A=2\sin a\cdot \sin b\cdot \cos (a+b)$}
\loigiai{
	Ta có
	\begin{eqnarray*}
	A&=&\sin^2(a+b)-\sin^2a-\sin^2b=\sin^2(a+b)-\dfrac{1-\cos 2a}{2}-\dfrac{1-\cos 2b}{2}\\
	&=&\sin^2(a+b)-1+\dfrac{1}{2}(\cos 2a+\cos 2b)=-\cos^2(a+b)+\cos (a+b)\cos (a-b)\\
	&=&\cos (a+b)\left[\cos (a-b)-\cos (a+b)\right]=2\sin a\sin b\cos (a+b).
	\end{eqnarray*}
}
\end{ex}

%%==========Câu 7
\begin{ex}%[1D1V3-1]
Giá trị lớn nhất của biểu thức $M=7\cos^2x-2\sin^2x$ là
\choice
{$-2$}
{$5$}
{\True $7$}
{$16$}
\loigiai{
	$M=7(1-\sin^2x)-2\sin^2x$ $=7-9\sin^2x$.\\
	Ta có $0\le\sin^2x\le 1$
	$\Leftrightarrow 0\ge -9\sin^2x\ge -9$, $\forall x\in \mathbb{R}$
	$\Leftrightarrow 7\ge 7-2\sin^2x\ge -2$.\\
	Giá trị lớn nhất là $7$.}
\end{ex}

%%==========Câu 8
\begin{ex}%[1D1H3-6]
Biết $A$, $B$, $C$ là các góc của tam giác $ABC$, khi đó.
\choice
{$\sin C=-\sin (A+B)$}
{$\cos C=\cos (A+B)$}
{$\tan C=\tan (A+B)$}
{\True $\cot C=-\cot (A+B)$}
\loigiai{
	Vì $A$, $B$, $C$ là các góc của tam giác $ABC$ nên $A+B+C=180^o\Rightarrow C=180^o-(A+B)$.\\
	Do đó $C$ và $(A+B)$ là 2 góc bù nhau.\\
	$\Rightarrow \sin C=\sin (A+B)$; $\cos C=-\cos (A+B)$; $\tan C=-\tan (A+B)$; $\cot C=-\cot (A+B)$.}
\end{ex}

%%==========Câu 9
\begin{ex}%[1D1H3-2]
Biết $\cot x=\dfrac{3}{4}$, $\cot y=\dfrac{1}{7}$, $x$, $y$ đều là góc dương, nhọn thì:
\choice
{$x+y=\dfrac{\pi}{4}$}
{$x+y=\dfrac{2\pi}{3}$}
{\True $x+y=\dfrac{3\pi}{4}$}
{$x+y=\dfrac{5\pi}{6}$}
\loigiai{
	$\cot x=\dfrac{3}{4}\Rightarrow \tan x=\dfrac{4}{3};\cot y=\dfrac{1}{7}\Rightarrow \tan y=7$.\\
	$\tan(x+y)=\dfrac{\tan x+\tan y}{1-\tan x\cdot\tan y}=-1\Rightarrow x+y=\dfrac{3\pi}{4}$ vì $x+y>0$.}
\end{ex}

%%==========Câu 10
\begin{ex}%[1D1H3-5]
Hãy xác định hệ thức \textbf{sai}
\choice
{$\sin x\cdot{\cos^3}x-\cos x\sin^3x=\dfrac{\sin 4x}{4}$}
{$\sin^4x+\cos^4x=\dfrac{3+\cos 4x}{4}$}
{\True $\dfrac{1+\sin x}{\cos x}=\cot \left(\dfrac{\pi}{4}+\dfrac{x}{2}\right)$}
{$\cot^2x+\tan^2x=\dfrac{2\cos 4x+6}{1-\cos 4x}$}
\loigiai{
	$\sin x\cdot \mathrm{co}\mathrm{s^3}x-\cos x\sin^3x=\sin x\cdot \cos x(\cos^2x-\sin^2x)=\dfrac{1}{2}\sin 2x\cdot \cos2x=\dfrac{\sin 4x}{4}$.\\
	$\sin^4x+\cos^4x=1-2\sin^2x\cdot \cos^2x=1-\dfrac{1}{2}\sin^22x=1-\dfrac{1}{2}(\dfrac{1-\cos 4x}{2})=\dfrac{3+\cos 4x}{4}$.\\
	$\dfrac{1+\sin x}{\cos x}=\dfrac{1-\cos\left(\dfrac{\pi}{2}+x\right)}{\sin\left(\dfrac{\pi}{2}+x\right)}=\dfrac{2\sin^2\left(\dfrac{\pi}{4}+\dfrac{x}{2}\right)}{2\sin\left(\dfrac{\pi}{2}+x\right)\cos\left(\dfrac{\pi}{4}+\dfrac{x}{2}\right)}=\tan \left(\dfrac{\pi}{4}+\dfrac{x}{2}\right)$.\\
	$\cot^2x+\tan^2x=\dfrac{\cos^2x}{\sin^2x}+\dfrac{\sin^2x}{\cos^2x}=\dfrac{\cos^4x+\sin^4x}{\cos^2x\cdot{\sin^2}x}=\dfrac{\dfrac{3+\cos 4x}{4}}{\dfrac{1-\cos 4x}{8}}=\dfrac{2\cos 4x+6}{1-\cos 4x}$.}
\end{ex}

%%==========Câu 11
\begin{ex}%[1D1V3-5]
	Nếu biết $\tan \alpha =\dfrac{1}{2}(0<a<90^\circ)$, $\tan b=-\dfrac{1}{3}(90^\circ<b<180^\circ)$ thì $\cos(2a-b)$ có giá trị đúng bằng
	\choice
	{\True $-\dfrac{\sqrt{10}}{10}$}
	{$\dfrac{\sqrt{10}}{10}$}
	{$-\dfrac{\sqrt{5}}{5}$}
	{$\dfrac{\sqrt{5}}{5}$}
	\loigiai{
		$\tan \alpha =\dfrac{1}{2}\Rightarrow \cos2a=\dfrac{1-\dfrac{1}{4}}{1+\dfrac{1}{4}}=\dfrac{3}{5}\Rightarrow \sin 2a=\dfrac{4}{5}$.\\
		$\tan b=-\dfrac{1}{3}(90^\circ<b<180^\circ)\Rightarrow \cos b=\dfrac{-4}{\sqrt{1+(-\dfrac{1}{3}})^2}=\dfrac{-3}{\sqrt{10}}$.\\
		$\Rightarrow \sin b=\tan b\cdot \cos b=\dfrac{-1}{3}\cdot \dfrac{-3}{\sqrt{10}}=\dfrac{1}{\sqrt{10}}$.\\
		$\Rightarrow \cos(2a-b)=\cos 2a\cos b+\sin 2a\sin b=\dfrac{3}{5}\cdot \dfrac{-3}{\sqrt{10}}+\dfrac{4}{5}\cdot \dfrac{1}{\sqrt{10}}=\dfrac{-1}{\sqrt{10}}$.}
\end{ex}

%%==========Câu 12
\begin{ex}%[1D1V3-1]
Giá trị nhỏ nhất của $M=\sin^4x+\cos^4x$ là
\choice
{$0$}
{$\dfrac{1}{4}$}
{\True $\dfrac{1}{2}$}
{$1$}
\loigiai{
	$M=\sin^4x+\cos^4x=1-\dfrac{1}{2}\sin^22x\ge 1-\dfrac{1}{2}=\dfrac{1}{2}$.\\
	Dấu bằng xảy ra khi $x=\dfrac{\pi}{4}+k\dfrac{\pi}{2}$, $k\in \mathbb{Z}$.}
\end{ex}

\Closesolutionfile{ans}
\indapan{6}{ans/ans-1-C3B3-DE2-LC}
\cauds
\Opensolutionfile{ans}[ans/ans-1-C3B3-DE2-DS]

%%==========Câu 1
\begin{ex}%[1D1H3-2]
	Biết $\sin a=\dfrac{8}{17}$, $\tan b=\dfrac{5}{12}$ và $a$, $b$ là các góc nhọn. Khi đó
	\choiceTF
	{\True $\tan a=\dfrac{8}{15}$}
	{\True $\sin (a-b)=\dfrac{21}{221}$}
	{$\cos (a+b)=\dfrac{14}{22}$}
	{$\tan (a+b)=\dfrac{17}{14}$}
	\loigiai{
		\begin{itemchoice}
			\itemch Đúng. Vì $a$, $b$ là các góc nhọn nên $\cos a>0$, $\cos b>0$.\\
			Ta có $\cos a=\sqrt{1-\sin ^2 a}=\dfrac{15}{17} \Rightarrow \tan a=\dfrac{\sin a}{\cos a}=\dfrac{8}{15}$.
			\itemch Đúng.  $\cos b=\sqrt{\dfrac{1}{1+\tan ^2 b}}=\dfrac{12}{13} \Rightarrow \sin b=\cos b \tan b=\dfrac{5}{13}$.\\ Khi đó $\sin (a-b)=\sin a \cos b-\cos a \sin b=\dfrac{8}{17} \cdot \dfrac{12}{13}-\dfrac{15}{17} \cdot \dfrac{5}{13}=\dfrac{21}{221}$.
			\itemch Sai. $\cos (a+b)=\cos a \cos b-\sin a \sin b=\dfrac{15}{17} \cdot \dfrac{12}{13}-\dfrac{8}{17} \cdot \dfrac{5}{13}=\dfrac{140}{221}$.
			\itemch Sai. $\tan (a+b)=\dfrac{\tan a+\tan b}{1-\tan a \tan b}=\dfrac{\dfrac{8}{15}+\dfrac{5}{12}}{1-\dfrac{8}{15} \cdot \dfrac{5}{12}}=\dfrac{171}{140}$.
		\end{itemchoice}
	}
\end{ex}

%%==========Câu 2
\begin{ex}%[1D1H3-2]
	Cho $\sin \alpha=\dfrac{2}{3}$, $\dfrac{\pi}{2}<\alpha<\pi$. Khi đó
	\choiceTF
	{\True $\cos \alpha=-\dfrac{\sqrt{5}}{3}$}
	{\True $\tan \alpha =-\dfrac{2\sqrt{5}}{5}$}
	{$\cos \left(\dfrac{\pi}{3}+\alpha\right)=\dfrac{\sqrt{5}-2\sqrt{3}}{6}$}
	{$\cos \left(\dfrac{\pi}{4}-\alpha\right)=\dfrac{\sqrt{10}-2\sqrt{2}}{6}$}
	\loigiai{
		\begin{itemchoice}
			\itemch Đúng. $\sin \alpha=\dfrac{2}{3}$, $\dfrac{\pi}{2}<\alpha<\pi$.\\
			Ta có $\cos ^2 \alpha=1-\sin ^2 \alpha=1-\left(\dfrac{2}{3}\right)^2=\dfrac{5}{9} \Rightarrow \cos \alpha= \pm \dfrac{\sqrt{5}}{3}$.\\
			Vì $\dfrac{\pi}{2}<\alpha<\pi$ nên $\cos \alpha=-\dfrac{\sqrt{5}}{3}$.
			\itemch Đúng. $\tan \alpha=\dfrac{\sin \alpha}{\cos \alpha}=-\dfrac{2 \sqrt{5}}{5}$.
			\itemch Sai. $\cos \left(\dfrac{\pi}{3}+\alpha\right)=\cos \dfrac{\pi}{3} \cos \alpha-\sin \dfrac{\pi}{3} \sin \alpha=\dfrac{1}{2} \cdot\left(\dfrac{-\sqrt{5}}{3}\right)-\dfrac{\sqrt{3}}{2} \cdot \dfrac{2}{3}=\dfrac{-\sqrt{5}-2 \sqrt{3}}{6}$.
			\itemch Sai. $\cos \left(\dfrac{\pi}{4}-\alpha\right)=\cos \dfrac{\pi}{4} \cos \alpha+\sin \dfrac{\pi}{4} \sin \alpha=\dfrac{\sqrt{2}}{2} \cdot\left(\dfrac{-\sqrt{5}}{3}\right)+\dfrac{\sqrt{2}}{2} \cdot \dfrac{2}{3}=\dfrac{-\sqrt{10}+2 \sqrt{2}}{6}$.
		\end{itemchoice}
	}
\end{ex}

%%==========Câu 3
\begin{ex}%[1D1H3-3]
	Biết $\cos 2 \alpha=\dfrac{5}{9}$, $0^\circ<\alpha<90^\circ$. Khi đó
	\choiceTF
	{\True $\sin \alpha=\dfrac{\sqrt{28}}{9}$}
	{\True $\cos \alpha=\dfrac{\sqrt{53}}{9}$}
	{$\tan \alpha =\dfrac{\sqrt{371}}{53}$}
	{\True $\cot \alpha =\dfrac{\sqrt{371}}{14}$}
	\loigiai{
		\begin{itemchoice}
			\itemch Đúng. Ta có $\cos 2 \alpha=\dfrac{5}{9}$, $0^\circ<\alpha<90^\circ$ nên 
			$\sin ^2 \alpha=\dfrac{1-\cos 2 \alpha}{2}=\dfrac{1-\left(\dfrac{5}{9}\right)^2}{2}=\dfrac{28}{81}\\ \Rightarrow \sin \alpha= \pm \dfrac{\sqrt{28}}{9}$.\\
			Vì $0^\circ<\alpha<90^\circ$ nên $\sin \alpha=\dfrac{\sqrt{28}}{9}$.
			\itemch Đúng.  $\cos ^2 \alpha=1-\sin ^2 \alpha=1-\left(\dfrac{\sqrt{28}}{9}\right)^2=\dfrac{53}{81} \Rightarrow \cos \alpha= \pm \dfrac{\sqrt{53}}{81}$.\\
			Vì $0^\circ<\alpha<90^\circ$ nên $\cos \alpha=\dfrac{\sqrt{53}}{9}$.
			\itemch Sai. $\tan \alpha=\dfrac{\sin \alpha}{\cos \alpha}=\dfrac{2 \sqrt{371}}{53}$.
			\itemch Đúng. $\cot \alpha=\dfrac{1}{\tan \alpha}=\dfrac{\sqrt{371}}{14}$.
		\end{itemchoice}
	}
\end{ex}

%%==========Câu 4
\begin{ex}%[1D1H3-3]
	Cho $\sin x=\dfrac{1}{5}$, $\dfrac{\pi}{2}<x<\pi$. Khi đó
	\choiceTF
	{$\sin 2x=\dfrac{4\sqrt{6}}{5}$}
	{\True $\cos 2x=\dfrac{23}{25}$}
	{$\tan 2x=\dfrac{20\sqrt{6}}{3}$}
	{$\cot 2x=\dfrac{23\sqrt{6}}{120}$}
	\loigiai{
		\begin{itemchoice}
			\itemch Sai. $\dfrac{\pi}{2}<x<\pi \Rightarrow \dfrac{2\pi}{2}<2x<2\pi \Rightarrow \pi <2x<2\pi \Rightarrow \sin 2x<0.\\
			\dfrac{\pi}{2}<x<\pi \Rightarrow \cos x<0 \\
			\sin x=\dfrac{1}{5}\Rightarrow \cos x=-\sqrt{1-\sin^2x}=-\dfrac{2\sqrt{6}}{5}.\\
			\sin 2x=2\sin x\cdot \cos x=2\cdot \dfrac{1}{5}\cdot \left(-\dfrac{2\sqrt{6}}{5}\right)=-\dfrac{4\sqrt{6}}{5}$.
			\itemch Đúng.  $\dfrac{\pi}{2}<x<\pi \Rightarrow \dfrac{2\pi}{2}<2x<2\pi \Rightarrow \pi <2x<2\pi \Rightarrow \cos 2x>0.\\
			\cos 2x=1-2\sin^2x=1-2\cdot \dfrac{1}{25}=\dfrac{23}{25}$.
			\itemch Sai. $\tan 2x=\dfrac{\sin 2x}{\cos 2x}=\dfrac{-\dfrac{4\sqrt{6}}{5}}{\dfrac{23}{25}}=-\dfrac{20\sqrt{6}}{3}$.
			\itemch Sai. $\cot 2x=\dfrac{\cos 2x}{\sin 2x}=\dfrac{\dfrac{23}{25}}{-\dfrac{4\sqrt{6}}{5}}=-\dfrac{23\sqrt{6}}{120}$.
		\end{itemchoice}
	}
\end{ex}
\Closesolutionfile{ans}
\indapan{3}{ans/ans-1-C3B3-DE2-DS}

\Opensolutionfile{ans}[ans/ans-1-C3B3-DE2-KQ]
\caukq

%%==========Câu 1
\begin{ex}%[1D1N3-2]
	Cho hai góc nhọn $a$ và $b$ với $\tan a=\dfrac{1}{7}$ và $\tan b=\dfrac{3}{4}$. Tính $a+b$.
	\shortans{$45$}
	\loigiai{
	Ta có $\tan (a+b)=\dfrac{\tan a+\tan b}{1-\tan a \tan b}=\dfrac{\dfrac{1}{7}+\dfrac{3}{4}}{1-\dfrac{1}{7} \cdot \dfrac{3}{4}}=1$, suy ra $a+b=45^\circ$.
	}
\end{ex}

%%==========Câu 2
\begin{ex}%[1D1V3-2]
	Cho các góc $\alpha$, $\beta$ thỏa mãn $\dfrac{\pi}{2}<\alpha$, $\beta<\pi$, $\sin \alpha=\dfrac{1}{3}$, $\cos \beta=-\dfrac{2}{3}$.\\
	Giá trị $\sin (\alpha +\beta)=-\dfrac{a+b\sqrt{c}}{d}$. Khi đó $bc-ad$ bằng bao nhiêu?
	\shortans{$2$}
	\loigiai{
	Do $\dfrac{\pi}{2}<\alpha$, $\beta<\pi \Rightarrow \heva{&\cos \alpha<0 \\& \sin \beta>0.}$\\
Ta có $\cos \alpha=-\sqrt{1-\sin ^2 \alpha}=-\sqrt{1-\dfrac{1}{9}}=-\dfrac{2 \sqrt{2}}{3};\\ \sin \beta=\sqrt{1-\cos ^2 \beta}=\sqrt{1-\dfrac{4}{9}}=\dfrac{\sqrt{5}}{3}$.\\
Suy ra $\sin (\alpha+\beta)=\sin \alpha \cdot \cos \beta+\cos \alpha \cdot \sin \beta=\dfrac{1}{3} \cdot\left(-\dfrac{2}{3}\right)+\left(-\dfrac{2 \sqrt{2}}{3}\right) \cdot \dfrac{\sqrt{5}}{3}=-\dfrac{2+2 \sqrt{10}}{9}$.\\
Do đó $\heva{&a=2 \\&b=2 \\&c=10 \\&d=9}\Rightarrow bc-ad=2$.	
	}
\end{ex}

%%==========Câu 3
\begin{ex}%[1D1H3-5]
	Cho $\cot \alpha=3$. Tính $A=\dfrac{3 \sin \alpha-2 \cos \alpha}{12 \sin ^3 \alpha+4 \cos ^3 \alpha}$ (kết quả làm tròn đến hàng phần mười).
	\shortans{$-0{,}3$}
	\loigiai{
	$A=\dfrac{3 \sin \alpha-2 \cos \alpha}{12 \sin ^3 \alpha+4 \cos ^3 \alpha}=\dfrac{\dfrac{1}{\sin ^2 \alpha}(3-2 \cot \alpha)}{12+4 \cot ^3 \alpha}=(1+\cot ^2 \alpha) \dfrac{3-2 \cot \alpha}{12+4 \cot ^3 \alpha}=-\dfrac{1}{4}\approx -0,3$.	
	}
\end{ex}

%%==========Câu 4
\begin{ex}%[1D1H3-6]
	Cho tam giác $ABC$ có $\tan B \sin ^2 C=\tan C \sin ^2 B$. Khi đó giá trị của $\widehat{B}-\widehat{C}$ bằng bao nhiêu độ?
	\shortans{$0$}
	\loigiai{
	}
\loigiai{
Ta có $0<\widehat{A}$, $\widehat{B}$, $\widehat{C}<180^\circ$.\\
Từ giả thiết, ta có $\dfrac{\sin B}{\cos B} \sin ^2 C=\dfrac{\sin C}{\cos C} \sin ^2B \Leftrightarrow \dfrac{1}{\cos B} \sin C=\dfrac{1}{\cos C} \sin B$\\ $\Leftrightarrow \sin C \cos C=\sin B \cos B \Leftrightarrow \dfrac{1}{2} \sin2C=\dfrac{1}{2} \sin ^2 B$\\
$\Leftrightarrow 2 \widehat{C}=2 \widehat{B}+k 2 \pi(k \in \mathbb{Z}) \Leftrightarrow 2 \widehat{C}=2 \widehat{B}$ (do $\widehat{B}$, $\widehat{C}$ là các góc của tam giác) $\Leftrightarrow \widehat{B}-\widehat{C}=0$.
	}
\end{ex}

%%==========Câu 5
\begin{ex}%[1D1V3-6]
	Cho tam giác $ABC$. Biết $\sin A+\sin B+\sin C=k\cos \dfrac{A}{2}\cos \dfrac{B}{2}\cos \dfrac{C}{2}$, tìm $k$?
	\shortans{$4$}
	\loigiai{Ta có
		\begin{eqnarray*}
		VT&=&\sin A+\sin B+\sin C \\
		& =&2 \sin \dfrac{A+B}{2} \cos \dfrac{A-B}{2}+2 \sin \dfrac{C}{2} \cos \dfrac{C}{2} \\
		& =&2 \sin \dfrac{\pi-C}{2} \cos \dfrac{A-B}{2}+2 \cos \dfrac{\pi-C}{2} \cos \dfrac{C}{2} \\
		& =&2 \cos \dfrac{C}{2} \cos \dfrac{A-B}{2}+2 \cos \dfrac{A+B}{2} \cos \dfrac{C}{2}\\
		& =&2\cos \dfrac{C}{2}\left(\cos \dfrac{A-B}{2}+\cos \dfrac{A+B}{2}\right) \\
		& =&2\cos \dfrac{C}{2}\cdot 2\cos \dfrac{A-B+A+B}{4}\cdot \cos \dfrac{A-B-A-B}{4} \\
		& =&4\cos \dfrac{C}{2}\cos \dfrac{A}{2}\cos \dfrac{-B}{2}\\
		&=&4\cos \dfrac{A}{2}\cos \dfrac{B}{2}\cos \dfrac{C}{2}.
		\end{eqnarray*}
	Vậy $\sin A+\sin B+\sin C=4 \cos \dfrac{A}{2} \cos \dfrac{B}{2} \cos \dfrac{C}{2}$.	
	}
\end{ex}

%%==========Câu 6
\begin{ex}%[1D1V3-7]
	Trong Vật lí, phương trình tổng quát của một vật dao động điều hoà cho bởi công thức $x(t)=A \cos (\omega t+\varphi)$, trong đó $t$ là thời điểm (tính bằng giây), $x(t)$ là li độ của vật tại thời điểm $t$, $A$ là biên độ dao động $(A>0)$ và $\varphi \in[-\pi; \pi]$ là pha ban đầu của dao động. Xét hai dao động điều hoà có phương trình: $x_1(t)=3 \cdot \cos \left(\dfrac{\pi}{6} t+\dfrac{\pi}{6}\right)(cm)$, $x_2(t)=3 \cdot \cos \left(\dfrac{\pi}{6} t+\dfrac{\pi}{4}\right)(cm)$. Pha ban đầu của dao động tổng hợp $x(t)=x_1(t)+x_2(t)$ là (đơn vị độ)
	\shortans{$37{,}5$}
	\loigiai{
	Ta có
	\begin{eqnarray*}
		x(t)&=&x_1(t)+x_2(t)=3\cos \left(\dfrac{\pi}{6}t+\dfrac{\pi}{6}\right)+3\cos \left(\dfrac{\pi}{6}t+\dfrac{\pi}{4}\right) \\
		& =&3\cdot2\left(\cos \left(\dfrac{\left. \dfrac{\pi}{6}t+\dfrac{\pi}{6}+\dfrac{\pi}{6}t+\dfrac{\pi}{4}\right)}{2}\right)\cos \left(\dfrac{\dfrac{\pi}{6}t+\dfrac{\pi}{6}-\dfrac{\pi}{6}t-\dfrac{\pi}{4}}{2}\right)\right) \\
		& =&6\cos \left(\dfrac{\pi}{6}t+\dfrac{5\pi}{24}\right)\cdot \cos \left(-\dfrac{\pi}{24}\right)=6\cos \left(\dfrac{\pi}{24}\right)\cdot \cos \left(\dfrac{\pi}{6}t+\dfrac{5\pi}{24}\right).
	\end{eqnarray*}
Vậy pha ban đầu của dao động tổng hợp là $\dfrac{5 \pi}{24}(\mathrm{rad})=37{,}5^\circ$.	
	}
\end{ex}
\Closesolutionfile{ans}
\indapan{6}{ans/ans-1-C3B3-DE2-KQ}